\chapter{주요 약어 및 용어}

{\footnotesize
\begin{longtable}{L{\dimexpr 0.1754\linewidth-2\tabcolsep\relax}L{\dimexpr 0.8246\linewidth-2\tabcolsep\relax}}
\caption{주요 약어 및 용어}\\
\toprule
\tabhead \textbf{약어} & \textbf{뜻} \\
\midrule
\endfirsthead
\toprule
\tabhead \textbf{약어} & \textbf{뜻} \\
\midrule
\endhead
\bottomrule
\endlastfoot
B\textsubscript{420}, B\textsubscript{470} & ISO 21498-1의 400 V급 부품 전압 하위 등급(상한 420 V, 470 V. 상한은 과도·리플을 제외한 정상 전압의 분류값). 이 보고는 사건 중에도 넘지 않는 설계 상한으로도 쓴다(3장) \\
CLLC & 변압기 양쪽에 인덕터와 커패시터 공진망을 두는 양방향 절연 공진형 변환기 \\
DAB & Dual-Active-Bridge. 변압기 양쪽의 능동 브리지를 이용하는 절연형 양방향 직류-직류 변환기 \\
DCX & DC Transformer. 공진점 고정 주파수로 운전해 이득이 1에 가까운 공진형 변환기(직류변압기) \\
EPS & Extended Phase Shift. DAB 한쪽 브리지의 두 레그 사이에도 위상차를 주는 변조(1차 내부 위상천이) \\
ODE & Ordinary Differential Equation. 시간에 따른 회로 상태를 표현하는 상미분방정식 \\
R2, R2-W410·W460 & 개선 구조와 그 창 조정 모듈(창 기준 410·460 V) \\
RSC & Resonant Switched-Capacitor Converter. 공진전류와 스위치드 커패시터를 이용하는 고정비 변환기 \\
SALS & Saturation-Aware Anticipatory Load Shaping. 조정단의 정격을 고려해 지연 가능한 부하를 승인·제한하는 상위 계층 \\
SPS & Single-Phase Shift. DAB 양쪽 브리지의 위상차로 전력을 조절하는 단일 위상천이 방식 \\
TB & Testbench. 검증 단계별로 남긴 Simulink 회로 모델 파일 \\
Tier-1 / Tier-2 & 계약 내 부하 변화 / 계약 위반·고장 포락선의 과도 시험 등급 \\
V\textsubscript{A}, V\textsubscript{W} & Stage A 출력·충전기 접속 노드의 전압 / 창 조정 기준전압 \\
V\textsubscript{A} 클램프 & 충전 중 배터리 경로를 잃었을 때 V\textsubscript{A} 노드의 충전기 전류를 저항으로 받는 보호 회로(최종 설계: 468/463 V, 가지 3개 × 2.4 Ω + 가지마다 SiC 스위치) \\
영교차 추종 게이트 & 위상마다 공진전류의 영교차를 검출해 다음 구간을 켜는 Stage A 게이트(검출 문턱 0.5 A) \\
A25, A25R & Stage A 25 kHz 설계(위치당 C4MS008075U2 2병렬, L\textsubscript{r} 864.6 nH, C\textsubscript{fly} 46.875 µF(해석값; 부품 후보 45.0 µF, 제어 타이밍은 45.0 µF 기준)). A25R은 V\textsubscript{A} 노드 커패시터를 2.4 mF로 키운 최종 설계 \\
C\textsubscript{A} & V\textsubscript{A} 노드 커패시터(최종 2.4 mF) \\
Q\textsubscript{oss} & 소자 출력용량에 V까지 충전된 전하. 하드 스위칭 전이 한 번의 손실은 Q\textsubscript{oss}(V)·V \\
P\textsubscript{rated} & 정격 충전점(375 A)의 위상 첨두 전류(A25R 74.5 A). 위상 첨두 판정의 기준 \\
M0·M1·D & SALS 부하 등급. M0는 감발하지 않는 부하(Stage C·V2L), M1은 한도를 줄일 수 있는 부하, D는 지연 가능한 부하 \\
F2 & 이중 고장의 Stage A 정지 판정(경로전류가 5 ms 안에 1 A 아래) \\
선정식 환산온도 & 클램프 스위치 판정 온도. 사건의 도통·스위칭 에너지를 첫 켬부터 마지막 끔까지의 시간으로 나눈 평균손실 × 선정용 정상 열저항 + 기준 온도(과도 열임피던스 계산이 아님) \\
켬·끔 지연(클램프) & 문턱 교차에서 클램프 스위치가 켜지거나 꺼질 때까지의 시간. 하드웨어 목표는 각각 ≤ 1 µs, 모의 판정의 입력은 각각 1.0 µs 고정(E44), 0.5 µs는 E45의 반례 확인 \\
ZCS / ZVS & Zero-Current / Zero-Voltage Switching. 영전류·영전압 스위칭 \\
공정 기준 & ④ 대 ⑤ 비교에서 ⑤의 보조 강압기에도 경부하 위상 셰딩을 준 조건(E41) \\
①–⑥ & ① 전용 승압기 + 800→400 V 강압기, ② Stage A + 전전력 조정기, ③ 기존 구성, ③′ 창 조정 변형(버스 ≤ 410 V), ④ 비조정 구조(버스 = V\textsubscript{A}, R2 기본형), ⑤ 모터·인버터 승압 + 800→400 V 강압기, ⑥ 모터·인버터 승압 + 800 V급 보조부품 \\
\end{longtable}
}
